% Part: proof-theory
% Chapter: natural-deduction
% Section: translation-G2i

\documentclass[../../../include/open-logic-section]{subfiles}

\begin{document}

\olfileid[pt]{pt}{nat}{tgi}

\olsection{Tradução de \Log{G2i} para \Log{N2i}}

Os antecedentes dos sequentes em \Log{G2i} são multiconjuntos~$\Gamma$
de !!{formula}s, enquanto os antecedentes dos sequentes em \Log{N2i}
são conjuntos~$\Gamma'$ de fórmulas rotuladas com a propriedade de que
cada rótulo ocorre apenas uma vez. Digamos que um conjunto~$\Gamma'$ de
!!{formula}s rotuladas \emph{corresponde a} um multiconjunto~$\Gamma$
se, para toda !!{formula}~$!A$, o número de !!{element}s de~$\Gamma'$ da
forma $x : !A$ é menor ou igual à multiplicidade de~$!A$ (isto é, ao
número de cópias de~$!A$) em~$\Gamma$.

\begin{prop}\ollabel{prop:G2i-to-N2i}
Se $\Log{G2i} + \Cut \Proves \Gamma \Sequent \Delta$, então $\Log{N2i}
\Proves \Gamma' \Sequent !C$, onde $!C \ident \lfalse$ se $\Delta$ é
vazio e $\Delta = \{!A\}$ caso contrário, e $\Gamma'$ corresponde
a~$\Gamma$.
\end{prop}

\begin{proof}
  Mostramos que, se $\pi$ é !!a{proof} de $\Gamma \Sequent \Delta$ em
  \Log{G2ci}, então existe um~$\Gamma'$ correspondente a~$\Gamma$ e
  !!a{proof}~$\delta$ de~$\Gamma' \Sequent !C$ em~\Log{N2i}. Fazemos
  isso por indução em $n = \pheight{\pi}$.

  Se $n = 0$, então $\pi$ é um axioma. Se for o axioma $\lfalse \Sequent
  \quad$, $\delta$ é o axioma $x:\lfalse \Sequent \lfalse$, isto é,
  $\Gamma' = \{x:\lfalse\}$ e $!C \ident \lfalse$. Se $\pi$ é um
  axioma da forma $!A \Sequent !A$, então $\delta$ é $x: !A \Sequent !A$.
  
  Se $n>0$, $\pi$ termina em uma regra~$R$. A hipótese de indução
  garante que há !!{proof}s de sequentes correspondentes às premissas
  dessa regra em~\Log{N2c}. Podemos supor que todos os rótulos de
  descarga nessas provas são distintos e que um rótulo $x$ ocorre
  apenas acima de uma inferência que o descarrega, caso seja de fato
  descarregado (renomeando-o, se necessário). Mostramos então como
  combinar essas !!{proof}s para obter provas de um sequente adequado
  $\Gamma' \Sequent !C$. Distinguimos casos conforme~$R$.

  \begin{enumerate}
    \item Se $R$ é \RightR{\Weakening}, $\pi$ termina em
    \[
    \AxiomC{}
    \RightLabel{$\pi_1$}
    \Deduce$\Gamma \fCenter$
    \RightLabel{\RightR{\Weakening}}
    \UnaryInf$\Gamma \fCenter !A$
    \DisplayProof
    \]
    A hipótese de indução aplicada a $\pi_1$ fornece !!a{proof}
    de~$\Gamma' \Sequent \lfalse$. Aplicamos $\FalseInt$ para obter
    !!a{proof} de $\Gamma' \Sequent !A$.
    \item Se $R$ é \LeftR{\Weakening}, $\pi$ termina em
    \[
    \AxiomC{}
    \RightLabel{$\pi_1$}
    \Deduce$\Gamma \fCenter \Delta$
    \RightLabel{\RightR{\Weakening}}
    \UnaryInf$!A, \Gamma \fCenter \Delta$
    \DisplayProof
    \]
    A hipótese de indução aplicada a $\pi_1$ fornece !!a{proof}
    de~$\Gamma' \Sequent !C$, que tomamos como $\delta$. Observe que, se
    $\Gamma'$ corresponde a $\Gamma$, também corresponde a $!A,
    \Gamma$, pois o número de $x : !A$ em $\Gamma'$ não é maior que o
    número de cópias de~$!A$ em $\Gamma$ e, portanto, também não é maior
    que o número de cópias de $!A$ em $\Gamma \cup \{!A\}$.
    \item Se $R$ é \LeftR{\Contraction}, $\pi$ termina em
    \[
    \AxiomC{}
    \RightLabel{$\pi_1$}
    \Deduce$!A, !A, \Gamma \fCenter \Delta$
    \RightLabel{\LeftR{\Contraction}}
    \UnaryInf$!A, \Gamma \fCenter \Delta$
    \DisplayProof
    \]
    A hipótese de indução aplicada a $\pi_1$ fornece !!a{proof}
    de~$\Pi, \Gamma' \Sequent !C$, onde $\Pi \subseteq \{x : !A,
    y:!A\}$. Se $\Pi = \{x : !A, y:!A\}$, substituímos todas as
    !!{formula}s rotuladas $y:!A$ em~$\pi_1$ por $x:!A$. Como $x$~ocorre
    no contexto do sequente final de~$\pi_1$, podemos supor que nenhuma
    inferência em~$\pi_1$ descarrega !!a{formula} da forma~$x:!B$, isto
    é, que nenhum $x:!B$ no lado esquerdo de um sequente em~$\pi_1$ é
    ativo. Consequentemente, o resultado ainda é !!a{proof} em~\Log{N2i}.
    \item Se $R$ é $\RightR{\lif}$, $\pi$ termina em
    \[
    \AxiomC{}
    \RightLabel{$\pi_1$}
    \Deduce$!A, \Gamma \fCenter !B$
    \RightLabel{\RightR{\lif}}
    \UnaryInf$\Gamma \fCenter !A \lif !B$
    \DisplayProof
    \]
    Pela hipótese de indução, existe !!a{proof}~$\delta_1$ de $x:
    !A, \Gamma' \Sequent !B$ ou $\Gamma' \Sequent !B$, onde $\Gamma'$
    corresponde a~$\Gamma$. Podemos tomar $\delta_1$ como $\delta$
    acrescida de uma aplicação de \Intro{\lif}.
    \item Se $R$ é \LeftR{\lif}, $\pi$ termina em
    \[
    \AxiomC{}
    \RightLabel{$\pi_1$}
    \Deduce$\Gamma_1 \fCenter !A$
    \AxiomC{}
    \RightLabel{$\pi_2$}
    \Deduce$!B, \Gamma_2 \fCenter \Delta$
    \RightLabel{\LeftR{\lif}}
    \BinaryInf$!A \lif !B, \Gamma_1, \Gamma_2 \fCenter \Delta$
    \DisplayProof
    \]
    A hipótese de indução fornece !!{proof}s $\delta_1$ de $\Gamma_1'
    \Sequent !A$ e $\delta_2$ de $x: !B, \Gamma_2' \Sequent !C$ ou de
    $\Gamma_2' \Sequent !C$. Se $\delta_2$ tem a segunda forma,
    $\delta_2$ pode servir como~$\delta$. Caso contrário, podemos
    renomear todas as fórmulas e todos os rótulos de descarga em~$\pi_1$
    para garantir que nenhum rótulo ocorra tanto em $\delta_1$ quanto
    em~$\delta_2$. Então $\Gamma_1' \cup \Gamma_2'$ corresponde a
    $\Gamma_1 \cup \Gamma_2$. Como $x:!B$ ocorre no sequente final
    de~$\delta_2$, $x:!B$ não pode estar ativo em nenhuma inferência com
    rótulo de descarga~$x$ em~$\delta_2$. Substituímos todo axioma $x:!B
    \Sequent !B$ em~$\delta_2$ pela !!{proof}~$\delta_3$
    \[
    \Axiom$x: !A \lif !B \fCenter !A \lif !B$
    \AxiomC{}
    \RightLabel{$\delta_1$}
    \Deduce$\Gamma_1' \fCenter !A$
    \RightLabel{\Elim{\lif}}
    \BinaryInf$x: !A \lif !B \fCenter !B$
    \DisplayProof
    \]
    e toda ocorrência de $x:!B$ em~$\delta_2$ por $x:!A \lif !B$.
    O resultado $\Subst{\delta_2}{\delta_3}{x:!B}$ é !!a{proof} de
    $x:!A \lif !B, \Gamma_1', \Gamma_2' \Sequent !C$.
    \item Se $R$ é \RightR{\land}, $\pi$ termina em
    \[
    \AxiomC{}
    \RightLabel{$\pi_1$}
    \Deduce$\Gamma_1 \fCenter !A$
    \AxiomC{}
    \RightLabel{$\pi_2$}
    \Deduce$\Gamma_2 \fCenter !B$
    \RightLabel{\RightR{\land}}
    \BinaryInf$\Gamma_1, \Gamma_2 \fCenter !A \land !B$
    \DisplayProof
    \]
    A hipótese de indução fornece uma prova $\delta_1$ de $\Gamma_1'
    \Sequent !A$ e $\delta_2$ de $\Gamma_2' \Sequent !A$. Renomeando
    todas as fórmulas e todos os rótulos de descarga em~$\pi_2$,
    podemos garantir que nenhum rótulo ocorra tanto em $\delta_1$ quanto
    em $\delta_2$. Então $\Gamma_1' \cup \Gamma_2'$ corresponde a
    $\Gamma_1 \cup \Gamma_2$. Obtemos !!a{proof}~$\delta$ aplicando
    $\Intro{\land}$ aos sequentes finais de $\delta_1$ e~$\delta_2$.
    \item Se $R$ é $\LeftR{\land}$, $\pi$ termina, por exemplo, em
    \[
    \AxiomC{}
    \RightLabel{$\pi_1$}
    \Deduce$!A, \Gamma \fCenter \Delta$
    \RightLabel{\LeftR{\land}}
    \UnaryInf$!A \land !B, \Gamma \fCenter \Delta$
    \DisplayProof
    \]
    Pela hipótese de indução, existe !!a{proof}~$\delta_1$ de $x:
    !A, \Gamma' \Sequent !C$ ou $\Gamma' \Sequent !C$, onde $\Gamma'$
    corresponde a~$\Gamma$. No segundo caso, podemos tomar $\delta_1$
    como $\delta$. No primeiro, observe que, como $x:!A$ ocorre no
    sequente final, $x:!A$ não está ativo em nenhuma inferência com
    rótulo de descarga~$x$ em~$\delta_1$. Substituímos todo axioma $x:!A
    \Sequent !A$ pela !!{proof}~$\delta_2$
    \[
    \Axiom$x: !A \land !B \fCenter !A \land !B$
    \RightLabel{\Elim{\land}}
    \UnaryInf$x: !A \land !B \fCenter !A$
    \DisplayProof
    \]
    e toda ocorrência de $x:!A$ por $x:!A \land !B$, isto é,
    consideramos $\Subst{\delta_1}{\delta_2}{x:!A}$. Essa é !!a{proof} de $x:!A
    \land !B, \Gamma' \Sequent !C$.
    \item Se $R$ é \Cut, $\pi$ termina em
    \[
    \AxiomC{}
    \RightLabel{$\pi_1$}
    \Deduce$\Gamma_1 \fCenter !A$
    \AxiomC{}
    \RightLabel{$\pi_2$}
    \Deduce$!A, \Gamma_2 \fCenter \Delta$
    \RightLabel{\Cut}
    \BinaryInf$\Gamma_1, \Gamma_2 \fCenter \Delta$
    \DisplayProof
    \]
    A hipótese de indução fornece uma prova $\delta_1$ de $\Gamma_1'
    \Sequent !A$ e $\delta_2$ de $x:!A, \Gamma_2' \Sequent !C$ ou de
    $\Gamma_2' \Sequent !C$. No segundo caso, tomamos $\delta$ como
    ~$\delta_2$. Caso contrário, substituímos todo axioma $x:!A \Sequent !A$ em
    $\delta_2$ por $\delta_1$ para obter~$\delta$ como a !!{proof}~$\Subst{\delta_2}{\delta_1}{x:!A}$ de $\Gamma_1',
    \Gamma_2' \Sequent !C$.
  \end{enumerate}
\end{proof}

\end{document}
